\documentclass[12pt]{article}
\usepackage{amsmath,amssymb,amsfonts}
\usepackage[margin=1in]{geometry}

\begin{document}

\section*{Problem 5.9}

\textbf{Problem:} Compute the three-dimensional Fourier transforms of the following functions:

\begin{itemize}
\item[(a)] $H(\mathbf{r}) = \dfrac{e^{-\alpha r}}{r}$
\item[(b)] $F(\mathbf{r}) = \dfrac{1}{r}$ \quad [To save work consider $F(r)$ as a special case of $H(r)$.]
\item[(c)] $F(\mathbf{r}) = e^{-r^2}$ \quad [The Fourier transform of this function can be obtained from that of $H(r)$ also.]
\end{itemize}

\bigskip
\textbf{Solution:}

The three-dimensional Fourier transform is:
\[
\tilde{f}(\mathbf{k}) = \frac{1}{(2\pi)^{3/2}}\int f(\mathbf{r})\,e^{i\mathbf{k}\cdot\mathbf{r}}\,d^3r.
\]

For spherically symmetric functions $f(\mathbf{r}) = f(r)$, choosing $\mathbf{k}$ along the $z$-axis:
\[
\tilde{f}(k) = \frac{1}{(2\pi)^{3/2}}\int_0^{\infty}\int_0^{\pi}\int_0^{2\pi} f(r)\,e^{ikr\cos\theta}\,r^2\sin\theta\,d\phi\,d\theta\,dr.
\]

Performing the $\phi$ and $\theta$ integrals:
\[
\tilde{f}(k) = \frac{1}{(2\pi)^{1/2}}\cdot\frac{2}{k}\int_0^{\infty} r\,f(r)\sin(kr)\,dr.
\]

\subsection*{Part (a): $H(r) = e^{-\alpha r}/r$}

\[
\tilde{H}(k) = \sqrt{\frac{2}{\pi}}\cdot\frac{1}{k}\int_0^{\infty} r\cdot\frac{e^{-\alpha r}}{r}\sin(kr)\,dr = \sqrt{\frac{2}{\pi}}\cdot\frac{1}{k}\int_0^{\infty} e^{-\alpha r}\sin(kr)\,dr.
\]

Using the standard integral $\int_0^{\infty}e^{-\alpha r}\sin(kr)\,dr = \frac{k}{\alpha^2 + k^2}$:
\[
\tilde{H}(k) = \sqrt{\frac{2}{\pi}}\cdot\frac{1}{k}\cdot\frac{k}{\alpha^2+k^2}.
\]

\[
\boxed{\tilde{H}(k) = \sqrt{\frac{2}{\pi}}\cdot\frac{1}{\alpha^2 + k^2}.}
\]

\subsection*{Part (b): $F(r) = 1/r$}

This is the special case $\alpha = 0$ of part (a):
\[
\boxed{\tilde{F}(k) = \sqrt{\frac{2}{\pi}}\cdot\frac{1}{k^2}.}
\]

This is the Fourier transform of the Coulomb potential, fundamental in electrostatics. The result shows that the Coulomb potential in $k$-space is $\propto 1/k^2$, which is Poisson's equation in Fourier space.

\subsection*{Part (c): $F(r) = e^{-r^2}$}

\[
\tilde{F}(k) = \sqrt{\frac{2}{\pi}}\cdot\frac{1}{k}\int_0^{\infty} r\,e^{-r^2}\sin(kr)\,dr.
\]

We evaluate $I = \int_0^{\infty}r\,e^{-r^2}\sin(kr)\,dr$. Using the identity:
\[
\int_0^{\infty}r\,e^{-r^2}\sin(kr)\,dr = \frac{k}{4}\sqrt{\pi}\,e^{-k^2/4},
\]
which can be derived by writing $\sin(kr) = \text{Im}(e^{ikr})$, completing the square, and using the Gaussian integral.

Proof: $\int_0^{\infty}r\,e^{-r^2}e^{ikr}\,dr = \int_0^{\infty}r\,e^{-(r-ik/2)^2-k^2/4}\,dr$. Taking the imaginary part and using the known result for the first moment of a shifted Gaussian gives the result above.

Therefore:
\[
\tilde{F}(k) = \sqrt{\frac{2}{\pi}}\cdot\frac{1}{k}\cdot\frac{k\sqrt{\pi}}{4}\,e^{-k^2/4} = \frac{1}{2\sqrt{2}}\,e^{-k^2/4}.
\]

\[
\boxed{\tilde{F}(k) = \frac{1}{2^{3/2}}\,e^{-k^2/4}.}
\]

A Gaussian transforms into a Gaussian---this is a fundamental property of the Fourier transform.

\end{document}
